% !TEX root = chapter-standalone.tex

\section{New from old}
\label{sec:new-from-old}

\todotext{J: @J: add material}



\subsection{Opposites}

\todotext{Add examples to concretely illustrate opposites. One idea is to use the list construction monoid as an example. Maybe some other examples are also very fitting and pedagogically nice.} 

\paragraph{Opposite semigroup}

Any semigroup has an opposite: it is obtained by ``flipping'' the order of composition.

\todotext{Rewrite definition so that 'in-line formula' is the main statement, and then the part about using the braiding function follows after as a remark.}

\begin{definition}[Opposite semigroup]\label{def:opposite-semigroup}
    $\sgrpA = \tup{\sgrpAset, {{\mtimes}}}$ be a \SY{semigroup}.
    Its \emph{opposite} $\sgrpA\op$ is the semigroup $\sgrpA\op = \tup{\sgrpAset, {{\mtimes \op}}}$ whose composition operation $\mtimes \op$ is the composite
    \begin{equation}
        \sgrpAset \cartprod \sgrpAset \overset{\braiding}{\fto} \sgrpAset \cartprod \sgrpAset \overset{\mtimes}{\fto} \sgrpAset
    \end{equation}
    where $\braiding$ is the braiding function
    \begin{equation}
        \defmapperiodset{
            \braiding
        }{
            \sgrpAset \cartprod \sgrpAset
        }{
            \sgrpAset \cartprod \sgrpAset
        }{
            \tup{\ela, \elb}
        }{
            \tup{\elb, \ela}
        }
    \end{equation}
    In other words,
    \begin{equation}
        \ela \mtimes \op \elb = \elb \mtimes \ela \qquad \forall \ela, \elb \setin \sgrpAset.
    \end{equation}

\end{definition}


\todotext{make the following a definition or combine somewhere with the other 'opposites' content; highlight connection to opposite semigroup notion, and that the identity element is the same.}

\begin{remark}
    Just like for semigroups, any monoid has an opposite.
    It is defined similarly: given a \SY{monoid} $\monoidAdefinition$, its opposite is $\monoidA \op = \tup{ \monoidAset, \mtimes\op, \idmon }$ where $\mtimes\op$ is defined by setting
    \begin{equation}
        \ela \mtimes \op \elb = \elb \mtimes_{\monoidA} \ela \qquad \forall \ela, \elb \setin \monoidAset.
    \end{equation}
\end{remark}

\subsection{Products}




\subsection{Sums}




\subsection{Sub-things}

\todotext{Not clear yet if we need/want to discuss 'subthings' at all. (One consideration: will be want to talk about subcategories?)}

\todotext{In 'subthings' examples, also have negative examples that show how things can fail.}

\paragraph{Subsemigroups}

\todotext{\bernina: Add some text and some examples to this subsection (subsemigroup)}

\begin{definition}[Subsemigroup]\label{def:subsemigroup}
    Let~$\sgrpA = \tup{\sgrpAset, \mtimes}$ be a \SY{semigroup}.
    A \maindef{subsemigroup} of~$\sgrpA$ is:

    \constit

    \begin{enumerate}
        \item A subset~$\sgrpBset \setsubseteq \sgrpAset$.
    \end{enumerate}

    \condit

    \begin{enumerate}
        \item The set~$\sgrpBset$ is closed under the composition operation~$\mtimes$ from~$\sgrpA$:
              \begin{equation}\label{eq:subsemigroup-closure}
                  \prfperiod{
                      \sgrpela \setin \sgrpBset \quad \sgrpelb \setin \sgrpBset
                  }{
                      (\sgrpela \mtimes \sgrpelb) \setin \sgrpBset
                  }
              \end{equation}
    \end{enumerate}
\end{definition}

\begin{example}
    \label{exa:evennum-subsemigroup}
    Consider the \SY{semigroup}~$\tup{\natnumbers,+}$ introduced in \cref{exa:natnum-semigroup}.
    We can take the subset of even natural numbers~$\sgrpBset \setsubset \natnumbers$.
    One can show that the sum of two even numbers is always even, satisfying the closure of~$\sgrpBset$ under the composition operation~$+$.
\end{example}



\paragraph{Submonoids}

\todotext{\bernina: Add some text and some examples to this subsection}

\begin{definition}[Submonoids]\label{def:submonoids}
    Let~$\monoidAdefinition$ be a \SY{monoid}.
    A \maindef{submonoid} of~$\monoidA$ is:

    \constit

    \begin{enumerate}
        \item A subset~$\monoidBset \setsubseteq \monoidAset$.
    \end{enumerate}

    \condit

    \begin{enumerate}
        \item The set~$\monoidBset$ is closed under the composition operation~$\mtimes$ from~$\monoidA$:
              \begin{equation}
                  \prfsemi{\monela \setin \monoidBset \quad \monelb \setin \monoidBset}{\monela \mtimes \monelb \setin \monoidBset}
              \end{equation}

        \item The \SY{neutral element}~$\idmon \setin \monoidAset$ is an element of~$\monoidBset$.
    \end{enumerate}
\end{definition}

\begin{example}
    $\tup{\natnumbers, +, 0}$,~$\tup{\wnumbers, +, 0}$, and~$\tup{\ratnumbers, +, 0}$ are all submonoids of~$\tup{\reals, +, 0}$.
\end{example}

\begin{example}
    $\tup{\natnumbers, \cdot, 1}$  is a \SY{submonoid} of~$\tup{\wnumbers, \cdot, 1}$.
\end{example}

\begin{example}
    $\makeset{ \ela \setin \natnumbers \mid \ela \text{ is even } }$ is a \SY{submonoid} of~$\tup{\natnumbers, +, 0}$.
\end{example}



\paragraph{Subgroups}

\todotextjira{454}{\bernina: Add more to subgroups subsection}

\begin{definition}[Subgroup]\label{def:subgroup}
    Let~$\grpA = \tup{\grpAset, \gtimes, \idgrp, \ginv}$ be a group.
    A \maindef{subgroup} of~$\grpA$ is:

    \constit

    \begin{enumerate}
        \item A subset~$\grpBset \setsubseteq \grpAset$.
    \end{enumerate}

    \condit

    \begin{enumerate}
        \item The set~$\grpBset$ is closed under the composition operation~$\gtimes$ from~$\grpA$:
              \begin{equation}
                  \prfsemi{
                      \grpela \setin \grpBset \quad \grpelb \setin \grpBset
                  }{
                      \grpela \gtimes \grpelb \setin \grpBset
                  }
              \end{equation}

        \item The \SY{neutral element}~$\idgrp \setin \grpAset$ is an element of~$\grpBset$;

        \item The set~$\grpBset$ is closed under the inverse operation~$\ginv$ from~$\grpA$:
              \begin{equation}
                  \prfperiod{
                      \grpela \setin \grpBset
                  }{
                      \ginv(\grpela) \setin \grpBset
                  }
              \end{equation}
    \end{enumerate}
\end{definition}





